\section{Introduction}

Large language models are increasingly consumed as remote frontier APIs, yet a growing community of users, researchers, and practitioners wants to \emph{self-host}: running open-weight models on hardware they own, for privacy, cost predictability, reproducibility, and independence from providers. The open question is not whether self-hosting is possible (mature serving stacks and quantized open weights make it routine) but whether it is \emph{viable}: does a self-hosted model deliver acceptable task accuracy at acceptable energy cost on the workloads people actually run? Energy is central to this question. The environmental footprint of ML is well documented at the training end \cite{schwartz2020green} and, more recently, at the inference end \cite{luccioni2024power}, and the software engineering community has called for empirical, measurement-driven approaches to Green AI \cite{cruz2024innovating}. However, published inference-energy figures overwhelmingly come from single-prompt benchmarks. They say little about \emph{agentic} workloads, in which an LLM is invoked hundreds of times inside a feedback loop, interleaved with other computation, and where a poorly steering model wastes its own inference joules and the joules of every downstream computation it triggers.

Agentic AutoML is a sharp instance of this workload class, and a deliberately chosen one: unlike open-ended agentic systems it has fixed-length inner experiments, a binary per-iteration outcome, and a ground-truth accuracy signal that is independent of the agent's own judgment. Those three properties are what make energy attributable to a phase and success measurable without a human oracle, most agentic workloads offer none of them. In Karpathy's \texttt{autoresearch} \cite{karpathy2026autoresearch}, an LLM \emph{proposer} iteratively mutates the training recipe of an inner machine-learning workload, launches a short training experiment for each mutation, inspects the outcome, and decides to keep or revert the change; in this study the inner workload is a compact convolutional network trained on CIFAR-10. A single session comprises tens of iterations, each combining LLM inference with a fixed, short burst of GPU training (45\,s here, set by calibration). The total energy cost of a session therefore depends jointly on (i) the proposer's architecture and serving configuration and (ii) the agent-loop configuration that governs how many experiments are run and how failures are handled. No published energy evaluation of \texttt{autoresearch}, or of any comparable agentic AutoML loop, exists to date.

This project evaluates a simple self-hostable model architecture under this workload. The architectural contrast is sparse versus dense: \texttt{Qwen3-30B-A3B}, a Mixture-of-Experts model holding 30.5B parameters but activating only $\sim$3B per token, against \texttt{Qwen3-14B}, a dense model from the same family, both quantized to int4 with AWQ and both fitting a single 20\,GB card. MoE sparsity is precisely the kind of ``simple architecture'' that could make self-hosting attractive, promising large-model quality at small-model inference cost; whether the theoretical per-token savings survive contact with a real agentic workload, and whether they cost final task accuracy, is an empirical question. Alongside the architectural factor we sweep two agent-loop configurations, \emph{patience} (how many consecutive regressions the agent tolerates before reverting) and \emph{loop budget} (the maximum number of inner experiments per session), each with a registered energy-relevance hypothesis.

Importantly, \texttt{autoresearch} is used here as a \emph{case study}, not as an object of scientific assessment: we evaluate its characteristics as a workload (single-machine, instrumentable, with a built-in ground-truth accuracy signal) and make no claims about the quality of the research it automates. The experiment is conducted according to the guidelines of Wohlin \etal \cite{wohlin12}, using the S2 group's \texttt{experiment-runner} framework \cite{karsten2025experimentrunner} as a session-level orchestrator, with \texttt{EnergiBridge} \cite{sallou2023energibridge} for energy measurement. The agent loop itself is re-implemented rather than adopted: the upstream repository publishes the protocol but no runnable loop, and, more importantly, the experiment requires loop policy to be \emph{enforced} rather than prompted and phase boundaries sharp enough to cut a power trace on (Section~\ref{sec:execution}). Because training and proposer inference are placed on physically separate GPUs, per-component energy is measured directly rather than estimated. A second experimental phase re-runs a subset of configurations with the proposer served on CPU only, for the self-hoster without a GPU.

From the results, practitioners will learn (i) whether an MoE architecture buys real energy savings on agentic workloads without losing accuracy, (ii) which agent-loop configurations sit on the energy/accuracy Pareto frontier, readable as direct configuration guidance (``at energy budget $X$, the most accurate configuration is $Y$''), and (iii) whether CPU-only proposer serving is a tolerable fallback. Researchers gain a reusable measurement pattern, session-level orchestration with per-device energy attribution, an error taxonomy that separates infrastructure failure from agent failure, and a keep threshold derived from a measured noise floor rather than assumed, together with a full replication package.
